% ch-a6-03.tex —— 第 6 篇（Feng--Harris）续：PDF 508--519 / 印刷 495--506
% 开头承接 ch-a6-02.tex 末尾中断的句子（4.3.2 节「奇异项」）。

\ysrc{495}情形 $a=0$。此时 $Z^r_{\mathcal E}(0)$ 按映射 $t_i\colon\mathcal E\boxtimes\mathcal O_S\to\mathcal F$ 的核分层，每个这样的层都贡献一项，其形式如上所述（一个 Chern 类乘以一个虚类，后者可像非奇异情形那样通过导出交构造）。数域上的类比情形有一个特殊性质：起支配作用的 Hermite 型在一般纤维上是正定的，这就排除了除一个层之外的所有层。更详细的讨论见 [FYZ21, §4] 与 [FYZ23, Introduction]。

结果是，对一切项 $a\in A_{\mathcal E}(\F_q)$，$[Z^r_{\mathcal E}(a)]^{\mathrm{vir}}$ 都能显式构造出来，但答案相当复杂。

另一条进路是：[FYZ21] 中观察到，特殊闭链 $Z^r_{\mathcal E}$ 的定义可以按 §3 的模式自然提升为导出特殊闭链（derived special cycles）$\mathcal Z^r_{\mathcal E}$：后者具有正确维数且拟光滑，其经典截断回到 $Z^r_{\mathcal E}$。据 [Kha19]，任何拟光滑导出叠 $\mathcal Z$ 都有内在的基本类概念 $[\mathcal Z]^{\mathrm{naive}}$。\footnote{这一粗略陈述长期以来是导出代数几何中的一条民间定理，例如已在 [Lur04] 中宣布；它对 Gromov--Witten 或 Donaldson--Thomas 型计数几何的重要性也早已为人所知。不过 Khan 对 $[\mathcal Z]^{\mathrm{vir}}$ 的构造是文献中出现的最一般者，[FYZ21] 用到了其一般性的全部范围。（原书此处作 $[\mathcal Z]^{\mathrm{vir}}$，而正文作 $[\mathcal Z]^{\mathrm{naive}}$，记号不一致；原书 p.495。）}[FYZ21, Theorem 6.5] 证明 $[\mathcal Z^r_{\mathcal E}(a)]^{\mathrm{naive}}$ 与先前（用初等方法）定义的 $[Z^r_{\mathcal E}(a)]^{\mathrm{vir}}$ 一致。

为勾画 $\mathcal Z^r_{\mathcal E}$ 的定义，关键点是认识到：$\bigsqcup_{\mathcal E}Z^r_{\mathcal E}$ 可以描述成一个纤维积（忽略有限 \'etale 扭子，解释见 [FYZ21, §3.6]）
\begin{equation}\tag{4.3.1}\label{eq:a6-03-1}
\begin{tikzcd}[column sep=large, row sep=large]
\bigsqcup_{\mathcal E}Z^r_{\mathcal E} \arrow[r] \arrow[d] & \mathrm{Hk}^r_{\mathrm M} \arrow[d] \\
\mathrm M \arrow[r, "{(\mathrm{Id},\mathrm{Frob})}"] & \mathrm M\times\mathrm M
\end{tikzcd}
\end{equation}
其中 $\mathrm M$ 是例 3.2.8 中 Hitchin 叠的酉推广，$\mathrm{Hk}^r_{\mathrm M}$ 是 $\mathrm M$ 的某个 Hecke 对应。现在，$\mathrm M$ 与 $\mathrm{Hk}^r_{\mathrm M}$ 都有自然的提升，成为（拟光滑的）导出叠 $\mathcal M$ 与 $\mathrm{Hk}^r_{\mathcal M}$；例如 $\mathcal M$ 是例 3.2.14 的酉推广。实际上，$\mathcal Z^r_{\mathcal E}$ 的拟光滑性要点在于映射 $\mathrm{Hk}^r_{\mathcal M}\to\mathcal M$ 是拟光滑的。然后，我们把 $\bigsqcup_{\mathcal E}\mathcal Z^r_{\mathcal E}$ 定义为导出纤维积
\[
\begin{tikzcd}[column sep=large, row sep=large]
\bigsqcup_{\mathcal E}\mathcal Z^r_{\mathcal E} \arrow[r] \arrow[d] & \mathrm{Hk}^r_{\mathcal M} \arrow[d] \\
\mathcal M \arrow[r, "{(\mathrm{Id},\mathrm{Frob})}"] & \mathcal M\times\mathcal M
\end{tikzcd}
\]
其经典截断由 (4.3.1) 回到 $Z^r_{\mathcal E}$。和前面一样，有映射 $\mathcal Z^r_{\mathcal E}\to A_{\mathcal E}(\F_q)$；我们定义 $\mathcal Z^r_{\mathcal E}(a)$ 为映射 $\mathcal Z^r_{\mathcal E}\to A_{\mathcal E}(\F_q)$ 在 $a$ 上的纤维。

\begin{rem}{4.3.1}\label{thm:a6-03-1}最近，Madapusi 在 [Mad23] 中定义了志村簇上的导出特殊闭链。这一情形的关键区别在于：与函数域情形不同，他的构造目前还没有配套的模空间论解释。事实上，正确的模问题应当用「整体 shtuka」来表述；这一概念在数域情形尚未定义，尽管 Scholze 等人近期的工作 [Sch18] 为其存在性给出了诱人的线索。
\end{rem}

\ysrc{496}\subsection{高阶 theta 级数}

设 $n\ge1$ 且 $n\ge m\ge1$。则可以把 $[Z^r_{\mathcal E}(a)]^{\mathrm{vir}}$ 组装成一个 Fourier 级数（Fourier 参数 $a\in A_{\mathcal E}(\F_q)$），[FYZ21] 称之为高阶 theta 级数（higher theta series）。

设 $m\in\Z_{\ge1}$。叠 $\mathrm{Bun}_{U^-(2m)}$ 参数化三元组 $(\mathcal G,h)$，其中 $\mathcal G$ 是 $X'$ 上秩 $2m$ 向量丛的族，$h\colon\mathcal G\xrightarrow{\sim}\sigma^*\mathcal G^*$ 是斜 Hermite 结构（即 $\sigma^*h^*=-h$）。

设 $\mathrm{Bun}_{\widetilde P_m}$ 是四元组 $(\mathcal G,h,\mathcal E)$ 的模叠，其中 $(\mathcal G,h)\in\mathrm{Bun}_{U^-(2m)}$，$\mathcal E\subset\mathcal G$ 是拉格朗日子丛（即 $\mathcal E$ 的秩为 $m$，且复合 $\mathcal E\subset\mathcal G\xrightarrow{h}\sigma^*\mathcal G^*\to\sigma^*\mathcal E^*$ 为零）。于是 $\widetilde P_m$ 对应 $U^-(2m)$ 的西格尔抛物子群。

设 $\mathcal E':=\mathcal G/\mathcal E$。$\mathcal G$ 上的斜 Hermite 型诱导完美配对 $\mathcal E\times\sigma^*\mathcal E'\to\mathcal O_{X'}$，它把 $\mathcal E'$ 与 $\sigma^*\mathcal E^*$ 等同起来。于是有短正合列
\[
0\to\mathcal E\to\mathcal G\to\sigma^*\mathcal E^*\to 0,
\]
并记 $e_{\mathcal G,\mathcal E}\in\Ext^1(\sigma^*\mathcal E^*,\mathcal E)$ 为其扩张类。由于 $\mathcal G$ 上的 Hermite 结构，这个扩张类落在子空间 $\Ext^1_{\mathrm{Herm}}(\sigma^*\mathcal E^*,\mathcal E)\subset\Ext^1(\sigma^*\mathcal E^*,\mathcal E)$ 中。在 $\Ext^1_{\mathrm{Herm}}(\sigma^*\mathcal E^*,\mathcal E)$ 与 $\Hom(\mathcal E,\sigma^*\mathcal E^\vee)=A_{\mathcal E}(\F_q)$ 之间有 Serre 对偶配对，记作 $\langle e_{\mathcal G,\mathcal E},a\rangle\in\F_q$。

取非平凡加性特征 $\psi\colon\F_q\to\overline{\Q}^\times$。高阶 theta 级数是 $(\mathcal G,h,\mathcal E)\in\mathrm{Bun}_{\widetilde P_m}(\F_q)$ 的函数，它给这样的三元组配置类
\begin{equation}\tag{4.4.1}\label{eq:a6-03-2}
\begin{aligned}
\Theta^r(\mathcal G,h,\mathcal E):={}&\chi(\det\mathcal E)q^{n(\deg\mathcal E-\deg\omega_X)/2}\sum_{a\in A_{\mathcal E}(\F_q)}\psi(\langle e_{\mathcal G,\mathcal E},a\rangle)[Z^r_{\mathcal E}(a)]^{\mathrm{vir}}\\
&\in\mathrm{CH}_{(n-m)r}(\mathrm{Sht}^r_{U(n)}),
\end{aligned}
\end{equation}
其中 $\chi\colon\mathrm{Pic}_{X'}(\F_q)\to\overline{\Q}^\times$ 是特征标，它在 $\mathrm{Pic}_X(\F_q)$ 上的限制等于二次特征标 $\mathrm{Pic}_X(\F_q)\to\{\pm1\}$ 的 $n$ 次幂，后者由类域论对应双覆盖 $X'/X$。

\subsection{模性猜想}

类比 Kudla 在 [Kud04] 中解释的猜想（其表述本身也部分仍是猜测），[FYZ21, Conjecture 4.16] 预言这些高阶 theta 级数是自守形式意义下的模性对象（modular）。具体地说，即：

\begin{conj}{4.5.1}\label{thm:a6-03-2}（模性猜想）§4.4 中的函数 $\Theta^r(\mathcal G,h,\mathcal E)$ 与拉格朗日子丛 $\mathcal E\subset\mathcal G$ 的选取无关。
\end{conj}

\begin{rem}{4.5.2}\label{thm:a6-03-3}「模性猜想」之所以得名，是因为它蕴含 $\Theta^r(\mathcal G,h,\mathcal E)$ 下降为 $\mathrm{Bun}_{U^-(2m)}(\F_q)$ 上的函数，而这个函数借助 Weil 均匀化可以解释为与 $U^-(2m)$ 相伴随的 adele 双陪集空间上的函数。于是它蕴含 $\Theta^r$ 定义了 $U^-(2m)$ 的自守形式，取值在 $\mathrm{CH}_*(\mathrm{Sht}^r_{U(n)})$ 中。
\end{rem}

当 $r=0$ 时，猜想 4.5.1 就是经典 theta 函数的模性，这早已知道。当 $r=1$ 时，它与正交与酉志村簇上算术 theta 级数的模性猜想平行，这方面近来进展很多，Chao Li 的文章 [Li25] 中有解释。但当 $r>1$ 时，这是经典理论中没有先例的新现象。

\ysrc{497}我们指出，$[Z^r_{\mathcal E}(a)]^{\mathrm{naive}}=[Z^r_{\mathcal E}(a)]^{\mathrm{vir}}$ 的导出构造已经使 [FYZ21] 能够证明某些预期性质，而这些性质用初等定义无法触及。对志村簇上特殊闭链的虚类而言，这不构成问题；差别在于函数域情形中奇异项的性质更复杂，如前所述。描述 $[Z^r_{\mathcal E}(a)]^{\mathrm{naive}}$ 的长处在于它更统一、更概念化，并且「只」依赖于底层的导出叠 $\mathcal Z^r_{\mathcal E}(a)$；而 $[Z^r_{\mathcal E}(a)]^{\mathrm{vir}}$ 不仅依赖于 $Z^r_{\mathcal E}(a)$，还依赖于它的某个辅助展示。因此，撰写 [FYZ21] 时人们觉得，导出视角对证明模性会更有用，尽管当时并不清楚如何实现。这一看法至少已部分得到 [FYZ23] 与 [FK24] 的证实：它们证明，高阶 theta 级数限制到 $\mathrm{Sht}^r_{U(n)}\to(X')^r$ 的一般纤维上之后是模性的。

的确，[FYZ23] 与 [FK24] 深深浸淫于导出代数几何的形式体系，尽管 $[Z^r_{\mathcal E}(a)]^{\mathrm{vir}}$ 的导出构造在证明中只是一个小环节。一般纤维上的模性来自 Fourier 分析新形式中的计算，其中包括 Borel--Moore 同调群上的所谓算术 Fourier 变换（当 $r=0$ 时它退化为通常的有限 Fourier 变换，此时关乎经典 theta 函数的模性），以及导出向量丛上（motivic）层的导出范畴上的导出 Fourier 变换。这些情形由迹构造的高阶版本联系起来，后者给出推广经典「层--函数对应」的「层--闭链对应」。哪怕勾画论证也超出本讨论的范围，我们只引用 [FYZ23, §1, §2] 中的概述。

\begin{rem}{4.5.3}\label{thm:a6-03-4}我们预期，要超越一般纤维上的模性，将需要 §9.4 中表述的迹猜想（Trace Conjecture）。它给出定义特殊闭链虚基本类的又一种途径。这一猜想也可以纯粹在经典代数几何内部表述。
\end{rem}

\section{导出伽罗瓦形变环}

本节概述 [GV18] 对导出伽罗瓦形变环的研究。为呈现「大图景」，我们将略去许多技术细节。渴望更多细节的读者可参阅论文 [GV18]；另一个资源是 Oberwolfach 工作小组报告 [Arb21]。阅读本节之前，读者应当熟悉伽罗瓦形变环的经典理论。

\subsection{概览}

为便于读者定位，我们先概述本节内容。

\emph{第一步。}首先我们说明如何构造导出整体无限制伽罗瓦形变环 $\mathcal R$，它表示例 3.2.13 中的函子 $\mathrm{M}^{\mathrm{Gal}}_G$。\footnote{这需要对 $\overline\rho$ 作一些假设；并且按例 3.2.13 的字面表述，该函子只有在 $G$ 为伴随型时才存在，不过通过计入行列式与中心，可按标准方式加以推广。}

\ysrc{498}\begin{rem}{5.1.1}\label{thm:a6-03-5}按 Mazur 的一条民间猜想，$\mathcal R$ 实际上应当同构于某个经典交换环，即其高阶同伦群为零。然而这一猜想远未解决；把 $\mathcal R$ 当作导出环来处理，使我们能为某些目的「绕过」它。事实上，我们所需的关于 $\mathcal R$ 的关键事实，是用伽罗瓦上同调对其切复形所作的描述；对导出环 $\mathcal R$，我们无条件地能算出这一描述，理由见 §A.19。而知道 $\pi_0(\mathcal R)$ 的切复形则等价于 Mazur 的猜想。
\end{rem}

\emph{第二步。}第一步的形变环 $\mathcal R$ 没有反映局部条件，而我们期待对应于自守形式的伽罗瓦表示满足这些局部条件。第二步是要求局部条件与自守形式所见到的条件「匹配」，从而把形变环切小。这些局部条件中最微妙的一点涉及 $p$ 处的分解群；在那里，局部条件由 $p$-进 Hodge 理论支配。

为此，人们考虑局部伽罗瓦群的类似导出形变函子，然后通过相对于整体形变函子取导出纤维积来施加局部条件。我们注意到，局部形变函子通常是离散的，因此正是「取导出纤维积」这一步引入了真正的导出结构。

\begin{rem}{5.1.2}\label{thm:a6-03-6}对 $p$ 处的局部条件，我们只考虑从整 $p$-进 Hodge 理论现代视角看也许是最简单的情形：Fontaine--Laffaille 范围内的晶体伽罗瓦形变函子。这应看作第一个测试情形；超越它似乎是推广时首先遇到的明显障碍。
\end{rem}

这一步得到一个对象 $\mathcal R^{\mathrm{crys}}$，称为导出晶体整体伽罗瓦形变环。在许多情形下它的虚拟维数为 $-\delta$，其中 $\delta>0$ 是「缺陷」（defect），因此它不可能是经典的。

我们「知道」$\mathcal R^{\mathrm{crys}}$ 的切复形，用伽罗瓦上同调描述；但其同伦群的结构先验地是神秘的。它将在下一步确定。

\emph{第三步。}作为 $R=T$ 定理的导出推广，人们期待 $\mathcal R^{\mathrm{crys}}$ 作用在局部对称空间的上同调上，并使该上同调至少在一般情形下是 $\mathcal R^{\mathrm{crys}}$ 上的「自由」模。这一作用在 [GV18] 中于同伦群层面构造出来，与 $\pi_*(\mathcal R^{\mathrm{crys}})$ 的确定同时进行，所用工具是 Taylor--Wiles 方法的 Calegari--Geraghty 修正。于是，$\pi_*(\mathcal R^{\mathrm{crys}})$ 的计算原则上是纯伽罗瓦论问题，却用到了与自守形式的对应。此外，这一计算还解释了局部对称空间上同调中重数的一个有趣数值模式；这个数值模式正是 [GV18] 与 [Ven19] 的原始动机之一。

\subsection{整体无限制导出形变环}

现在回到开头，讨论如何构造 §5.1 第一步中的导出伽罗瓦形变函子。

\subsubsection{经典形变环}

设 $k$ 是有限域。暂时设 $\Gamma$ 是离散群，$G$ 是 Witt 向量 $W(k)$ 上的分裂代数群。

设 $\overline\rho$ 是 $\Gamma$ 在 $G(k)$ 中的表示。设 $\mathrm M^\Gamma_G$ 是范畴上的函子，该范畴的对象是在 $k$ 上带增广的 Artin 局部 $W(k)$-代数 $A$；函子把 $A$ 送到提升集合

\ysrc{499}模去共轭，
\begin{equation}\tag{5.2.1}\label{eq:a6-03-3}
\left\{
\begin{tikzcd}[column sep=large, row sep=large]
& G(A)\arrow[d]\\
\Gamma\arrow[r,"\overline\rho"']\arrow[ur,"\rho",dashed] & G(k)
\end{tikzcd}
\right\}/\sim
\end{equation}

若 $\overline\rho$ 绝对不可约，且 $\Gamma$ 的群上同调满足一定的有限性条件，则 Schlessinger 判别法（Schlessinger's criterion）蕴含 $\mathrm M^\Gamma_G$ 可表。这些观察，以及考虑伽罗瓦形变模空间这一原始想法，都归功于 Mazur [Maz89]。把 $\mathrm M^\Gamma_G$ 看作由 $\Gamma$ 的有限商所得诸形变函子的（滤过）余极限，就可以把这一故事推广到投射有限群 $\Gamma$。

设 $\mathcal O_F$ 是整体域 $F$ 的整数环，$S$ 是 $F$ 的有限位点集。此时 $\Gamma=\pi_1(\Spec\mathcal O_F[1/S])$ 由类域论满足所需的有限性条件。在这种情形下，我们实际上要把 $\mathrm M^\Gamma_G$ 记作 $\mathrm M^{\overline\rho}_{\mathcal O_F[1/S]}$，因为这一记号强调了本节将要关注的参数。其万有代表环是（经典的）整体无限制伽罗瓦形变环，记作 $\mathrm R^{\overline\rho}_{\mathcal O_F[1/S]}$。

\subsubsection{Galatius--Venkatesh 的工作}

Galatius--Venkatesh 的想法是：通过在单纯层面重复这个故事，构造伽罗瓦形变环的导出（derived）版本。论文 [GV18] 用的是单纯交换环与模型范畴的语言，而不是动态交换环与 $\infty$-范畴的语言。我们借此机会讨论这一对照所暴露的一些微妙之处。

从道德上说，我们想把 $\mathrm M^{\overline\rho}_{\mathcal O_F[1/S]}$ 提升为函子 $\mathcal M^{\overline\rho}_{\mathcal O_F[1/S]}$，定义在 $k$ 上带增广的动态 Artin 局部 $W(k)$-代数 $A$ 的范畴上（该范畴尚待适当定义）；它把这样的 $A$「送到」模掉共轭的提升构成的 anima：
\begin{equation}\tag{5.2.2}\label{eq:a6-03-4}
\left\{
\begin{tikzcd}[column sep=large, row sep=large]
& G(\mathcal A)\arrow[d]\\
\Gamma\arrow[r,"\overline\rho"']\arrow[ur,"\rho",dashed] & G(k)
\end{tikzcd}
\right\}/\sim
\end{equation}

由于 $\infty$-范畴之间函子的定义十分繁复——至少在 [Lur09] 的形式体系中如此——严格地构造函子可能很困难。特别地，$\infty$-范畴的函子不能靠指定对象与态射分别被送到哪里来给定（文献中有时对此处理得不够严谨），因此上一段的讨论其实并没有定义一个合法的函子。

Galatius--Venkatesh 选择用单纯交换环的模型范畴来应对这一困难。可以把它看作动态交换环的 $\infty$-范畴的一种刚性化（rigidification）。它是普通的 1-范畴，因此函子可按通常方式在对象与态射上指定。然而，模型范畴的结构中包含「弱等价」的概念，别名「同伦等价」；而我们只关心把同伦等价送到同伦等价的函子。这样的函子称为同伦不变的（homotopy invariant）。

\ysrc{500}为把它与另一视角联系起来：动态交换环的 $\infty$-范畴乃是单纯交换环在同伦等价处的局部化，因此从单纯交换环出发的同伦不变函子应当等价于从动态交换环出发的函子。这类陈述可以讲得更精确，但 [GV18] 的进路是自始至终使用模型论框架。

\subsubsection{导出形变环}

现在概述 $\mathrm R^{\overline\rho}_{\mathcal O_F[1/S]}$ 的构造。我们从形变函子 $\mathrm M^{\overline\rho}_{\mathcal O_F[1/S]}$ 的导出开始。首先需要定义 Artin 局部单纯交换环这一概念。

\begin{defn}{5.2.1}\label{thm:a6-03-7}单纯交换环 $A$ 称为 \emph{Artin 局部的}（Artinian local），若 $\pi_0(A)$ 是 Artin 局部的，且 $\pi_*(A)$ 作为 $\pi_0(A)$-模是有限生成的。
\end{defn}

接下来，我们粗略地想把 $\mathrm M^{\overline\rho}_{\mathcal O_F[1/S]}(A)$ 定义为把 $\overline\rho\colon\Gamma\to G(k)$ 提升为 $\rho\colon\Gamma\to G(A)$ 的单纯集，并模去等价。我们需要定义这个单纯集，并保证其构造以同伦不变的方式依赖于 $A$。为此，换一种关于表示的视角会很有用。

给定离散群 $H$，有分类空间 $BH$。分类空间 $BH$ 可定义为单纯集 $N_\bullet(H)$ 的几何实现，而 $N_\bullet(H)$ 是把 $H$ 看作范畴（只有一个对象，自同态由 $H$ 给出）所得的脉（nerve）。单纯集 $N_\bullet(H)$ 形如
\begin{equation}\tag{5.2.3}\label{eq:a6-03-5}
\begin{tikzcd}[row sep=small, column sep=small]
\ast & H\arrow[l, shift left=1.5pt]\arrow[l, shift right=1.5pt] & H\times H\arrow[l, shift left=3pt]\arrow[l]\arrow[l, shift right=3pt] & \cdots
\end{tikzcd}
\end{equation}
其中面映射与退化映射的显式公式可在 [May99, p.128] 中找到。特别地，$BH$ 带有一个区分出的基点 $\mathrm{pt}$。

若 $X$ 是「好」空间（例如某个单纯集的几何实现），则从 $X$ 到 $BH$ 的映射的同伦类与 $X$ 上的「$H$-局部系统」一一对应。若 $X$ 连通，则选定基点 $x\in X$ 后，这些又与同态 $\pi_1(X,x)\to H$ 模共轭的集合一一对应。因此，从 $X$ 到 $BH$ 的映射空间给出「$X$ 的局部系统空间」。（更多讨论见 [Arb21, C.4]。）

更一般地，公式 (5.2.3) 可用来把分类空间的定义推广到非离散的群。若 $H$ 是单纯群，则 (5.2.3) 是双单纯集，我们把单纯集 $N_\bullet(H)$ 定义为该双单纯集的对角线。这一推广在定义导出群表示函子时很重要：那里 $H$ 会通过把代数群 $G$ 作用在单纯交换环 $R$ 上而得到——正是从离散的 $R$ 过渡到单纯交换的 $R$，在 $H$ 上引入了非平凡的单纯结构。

\begin{exam}{5.2.2}\label{thm:a6-03-8}设 $\Gamma$ 是离散群，并转换视角，把 $B\Gamma$ 与 $BH$ 都视为单纯集。则由一般范畴论考虑，存在自然的单纯集，即态射 $B\Gamma\to BH$，应把它想成 $\Gamma$ 在 $H$ 中表示的单纯集。注意这一形式化已经包含了「模掉共轭」。若要刻画的是同态\tn{原文在此处作「homorphisms」，疑为「homomorphisms」之误。（原书 p.500）}而非表示，则取带基点的态射 $(B\Gamma,\mathrm{pt})\to(BH,\mathrm{pt})$。（依据是 [Lur17, §5.2.6.10, 5.2.6.13]：该处说明，从单纯群到带基点空间的分类空间函子是全忠实的。）
\end{exam}

\ysrc{501}至此，我们为「表示 $\Gamma\to G(A)$ 的单纯集」（模共轭）定义了一个自然候选：它应当是态射 $B\Gamma\to BG(A)$ 的单纯集，其中若 $\Gamma\cong\varprojlim_i\Gamma_i$ 投射有限，则把 $B\Gamma$ 视为诸 $B\Gamma_i$ 的形式反向极限。

[GV18] 的实际进路略有不同。为解释它，我们回忆：按 Artin--Mazur [AM86] 与 Friedlander [Fri82]，可以从概形 $X$ 提取出 \'etale 同伦型（\'etale homotopy type）$\text{\'Et}(X)$。这是 pro-单纯集，其用途是捕捉 $X$ 的 \'etale 拓扑。当 $\mathcal O_F[1/S]$ 是整体域 $F$ 中的 $S$-整数环时，$\text{\'Et}(\mathcal O_F[1/S])\cong B\Gamma$，其中 $\Gamma$ 是 $\Spec\mathcal O_F[1/S]$ 关于某个基点的 \'etale 基本群。（于是 $\Spec\mathcal O_F[1/S]$ 是一个「\'etale $K(\pi,1)$」，至少在 $p>2$ 的 $p$-进系数下如此。）我们更愿意用 $\text{\'Et}(\mathcal O_F[1/S])$ 而非 $B\Gamma$，尽管两者等价，因为前者不涉及基点的选取。\footnote{这与 §3 中关于 Betti 局部系统导出模空间的讨论类似：那里 $\Gamma$ 类比于 $\pi_1(C)$，而 $\text{\'Et}(\mathcal O_F[1/S])$ 类比于 $C$ 本身。}由前述讨论，$\mathrm M^\Gamma_G(A)$ 的模空间描述可重新表述为提升的集合
\[
\begin{tikzcd}[column sep=large, row sep=large]
& BG(A)\arrow[d]\\
\text{\'Et}(\mathcal O_F[1/S])\arrow[r,"\overline\rho"']\arrow[ur,"\rho",dashed] & BG(k)
\end{tikzcd}
\]

于是我们把 $\mathcal M^{\overline\rho}_{\mathcal O_F[1/S]}$ 定义为：在带 $k$ 上增广的 Artin 局部单纯交换环 $A$ 的范畴上的函子，它把 $A$ 送到「提升」的单纯集
\[
\begin{tikzcd}[column sep=large, row sep=large]
& BG(c(A))\arrow[d]\\
\text{\'Et}(\mathcal O_F[1/S])\arrow[r,"\overline\rho"']\arrow[ur,"\rho",dashed] & BG(k)
\end{tikzcd}
\]
其中 $c(A)$ 是 $A$ 的函子性余纤维替换。这一「余纤维替换」概念是模型范畴结构的另一部分，类似于同调代数中的投射消解；正是它保证该函子确实同伦不变。\footnote{一个微妙之处是：$c(-)$ 不是幺半的，因此 $BG(c(A))$ 指的是余单纯交换环 $W(k)\rightrightarrows c(\mathcal O_G)\mathrel{\substack{\longrightarrow\\[-2pt]\longrightarrow\\[-2pt]\longrightarrow}}c(\mathcal O_G\otimes\mathcal O_G)\ldots$，其中诸映射与 (5.2.3) 中的映射对偶。}这里「提升」按定义是指同伦纤维积
\[
\Map\bigl(\text{\'Et}(\mathcal O_F[1/S]),BG(c(A))\bigr)\times^{h}_{\Map(\text{\'Et}(\mathcal O_F[1/S]),BG(k))}\{\overline\rho\},
\]
它是在单纯集的模型范畴中先取纤维替换、再取纤维积来计算的。

\subsubsection{可表性}

Lurie 的博士论文 [Lur04] 建立了 Schlessinger 可表性判别法的导出版本。要检验该判别法，需要计算切复形；答案已在例 3.2.13 中提到，后文还会回顾。

\ysrc{502}利用 Lurie 的导出 Schlessinger 判别法，Galatius--Venkatesh 证明 $\mathcal M^{\overline\rho}_{\mathcal O_F[1/S]}$ 可由一个 pro Artin 动态环 $\mathcal R^{\overline\rho}_{\mathcal O_F[1/S]}$ 表示。在我们看来，检验该判别法并不是 [GV18] 中最有趣的部分，因此完全跳过。

\subsection{Desiderata（期望事项）}

我们把 $\mathcal M^{\overline\rho}_{\mathcal O_F[1/S]}$ 与 $\mathcal R^{\overline\rho}_{\mathcal O_F[1/S]}$ 的性质抽象出来，这些性质是后文真正会用到的。

\subsubsection{切复形}

如例 3.2.13 中已经讨论过的，$\mathcal M^{\overline\rho}_{\mathcal O_F[1/S]}$ 在 $\overline\rho$ 处的切复形是 $R\Gamma(\Spec\mathcal O_F[S^{-1}],\frakg_{\overline\rho}[1])$。这一点在 [GV18, Example 4.38] 中有解释；注意它与 §3 中讨论的 Betti 局部系统导出模空间十分相似。

\subsubsection{与经典形变函子的比较}

我们与经典理论有如下相容性。

作为 §2.3 的一个变体，存在从在 $k$ 上带增广的 Artin 局部环 $A$ 的范畴到在 $k$ 上带增广的 Artin 局部动态环的范畴的全忠实嵌入。由构造，我们有交换图
\[
\begin{tikzcd}[column sep=large, row sep=large]
\left\{\substack{\text{在 $k$ 上带增广的 Artin}\\\text{局部交换环 $A$}}\right\}\arrow[d]\arrow[r,"{\mathrm M^{\overline\rho}_{\mathcal O_F[1/S]}}"] & \mathsf{Set}\\
\left\{\substack{\text{在 $k$ 上带增广的 Artin}\\\text{局部单纯交换环 $\mathcal A$}}\right\}\arrow[r,"{\mathcal M^{\overline\rho}_{\mathcal O_F[1/S]}}"'] & \mathsf{sSet}\arrow[u,"\pi_0"']
\end{tikzcd}
\]
在代表环的层面，该图的交换性等价于同构
\[
\pi_0\bigl(\mathcal R^{\overline\rho}_{\mathcal O_F[1/S]}\bigr)\cong \mathrm R^{\overline\rho}_{\mathcal O_F[1/S]}.
\]

\begin{rem}{5.3.1}\label{thm:a6-03-9}如前面提到的，实际上有一条归于 Mazur 的民间猜想：$\mathrm R^{\overline\rho}_{\mathcal O_F[1/S]}$ 是维数为 $\chi(R\Gamma(\mathcal O_F[1/S],\frakg_{\overline\rho}[1]))$ 的 LCI。按 §3.2.1，这等价于说 $\mathcal R^{\overline\rho}_{\mathcal O_F[1/S]}$ 实际上同伦离散，即对 $i>0$ 有 $\pi_i(\mathcal R^{\overline\rho}_{\mathcal O_F[1/S]})=0$。

那么，既然本来可以直接用 $\mathrm R^{\overline\rho}_{\mathcal O_F[1/S]}$，为何还要费力定义 $\mathcal R^{\overline\rho}_{\mathcal O_F[1/S]}$？往后我们会用到的最重要事实，是关于 $\mathcal R^{\overline\rho}_{\mathcal O_F[1/S]}$ 的切复形的描述，而它来得相对容易。若不知道 Mazur 的猜想，我们对 $\mathrm R^{\overline\rho}_{\mathcal O_F[1/S]}$ 的切复形除了它的第 0 上同调群之外一无所知。
\end{rem}

\subsection{施加局部条件}

现在我们开始考虑施加局部条件。我们将从两方面用到它：
\begin{enumerate}[leftmargin=1.8em]
\item 施加 $p$-进 Hodge 理论型的条件，这是把范围切小到可能「匹配」自守侧的伽罗瓦表示所必需的。
\item 实现 Taylor--Wiles 方法，其中需要考虑在辅助素数处增加与去掉分歧的效应。
\end{enumerate}

\ysrc{503}\subsubsection{局部形变函子}

对局部域 $F_v$ 与 $\pi_1(F_v)$ 的表示 $\overline\rho$，可以类似整体情形定义局部（无限制）导出形变函子 $\mathcal M^{\overline\rho}_{F_v}$，它把
\[
A\mapsto \Map\bigl(\text{\'Et}(F_v),BG(c(A))\bigr)\times^{h}_{\Map(\text{\'Et}(F_v),BG(k))}\{\overline\rho\}.
\]
若 $v$ 使得 $\overline\rho|_{\pi_1(F_v)}$ 无分歧，则还有一个有用的变体 $\mathcal M^{\overline\rho}_{\mathcal O_{F_v}}$，它参数化无分歧形变，把
\[
A\mapsto \Map\bigl(\text{\'Et}(\mathcal O_{F_v}),BG(c(A))\bigr)\times^{h}_{\Map(\text{\'Et}(\mathcal O_{F_v}),BG(k))}\{\overline\rho\}.
\]
这些函子出现时通常不可表，因为 $\overline\rho$ 通常不可约。

\subsubsection{切出局部条件}

设 $v$ 是 $F$ 的一个位点。沿 $\Spec F_v\to\Spec\mathcal O_F[1/S]$ 拉回局部系统，得到一个映射
\[
\mathcal M^{\overline\rho}_{\mathcal O_F[1/S]}\to \mathcal M^{\overline\rho}_{F_v}.
\]
若 $v\notin S$，则上述映射自然地经过 $\mathcal M^{\overline\rho}_{\mathcal O_F[1/S]}\to\mathcal M^{\overline\rho}_{\mathcal O_{F_v}}$ 分解。

现在假设给定某个函子 $\mathcal M^{\overline\rho,D_v}_{F_v}\to\mathcal M^{\overline\rho}_{F_v}$，其中 $D$ 代表某个形变问题。通常 $\mathcal M^{\overline\rho,D_v}_{F_v}\to\mathcal M^{\overline\rho}_{F_v}$ 应当是闭浸入，意思是在经典截断上是闭浸入。一个例子是 $\mathcal M^{\overline\rho}_{\mathcal O_{F_v}}\hookrightarrow\mathcal M^{\overline\rho}_{F_v}$；非正式地说，它切出的正是「在 $v$ 处无分歧」的那部分形变。于是可以取下图的导出纤维积
\[
\begin{tikzcd}[column sep=large, row sep=large]
\mathcal M^{\overline\rho,\{D_v\}}_{\mathcal O_F[1/S]}\arrow[r]\arrow[d] & \mathcal M^{\overline\rho}_{\mathcal O_F[1/S]}\arrow[d]\\
\mathcal M^{\overline\rho,D_v}_{F_v}\arrow[r] & \mathcal M^{\overline\rho}_{F_v}
\end{tikzcd}
\]

\begin{exam}{5.4.1}\label{thm:a6-03-10}（增加分歧）设 $v\notin S$，令 $S'=S\sqcup\{v\}$。态射 $\Spec\mathcal O_F[1/S']\to\Spec\mathcal O_F[1/S]$ 诱导映射 $\mathcal M^{\overline\rho}_{\mathcal O_F[1/S]}\to\mathcal M^{\overline\rho}_{\mathcal O_F[1/S']}$。非正式地说，这个态射把一个在 $S$ 外无分歧的形变送到同一个形变，但把它看作在 $S'$ 外无分歧（即忘掉它在 $v$ 处无分歧这一事实）。

在经典形变环的理论中有
\[
\mathrm M^{\overline\rho}_{\mathcal O_F[1/S']}\times_{\mathrm M^{\overline\rho}_{F_v}}\mathrm M^{\overline\rho}_{\mathcal O_{F_v}}\cong\mathrm M^{\overline\rho}_{\mathcal O_F[1/S]}.
\]
非正式地说，这等价于如下陈述：「$\mathcal O_F[1/S']$ 上在 $v$ 处无分歧的局部系统的形变，唯一地来自 $\mathcal O_F[1/S]$ 上局部系统的形变」，而这是显然的。然而我们特意把这条陈述单列出来，理由有二：(i) 它在 Taylor--Wiles 方法中从粘合形变环下降时起关键作用；(ii) 它的导出版本并不平凡。

为回应第 (ii) 点：可以证明类似的等式在导出形变函子层面也成立：
\begin{equation}\tag{5.4.1}\label{eq:a6-03-6}
\mathcal M^{\overline\rho}_{\mathcal O_F[1/S']}\times^{h}_{\mathcal M^{\overline\rho}_{F_v}}\mathcal M^{\overline\rho}_{\mathcal O_{F_v}}\cong\mathcal M^{\overline\rho}_{\mathcal O_F[1/S]}.
\end{equation}

\ysrc{504}这由 [GV18] 通过计算切复形的映射来证明。粗略地说，(5.4.1) 的内容是：下图
\[
\begin{tikzcd}[column sep=large, row sep=large]
\Spec F_v\arrow[r]\arrow[d] & \Spec\mathcal O_F[1/S']\arrow[d]\\
\Spec\mathcal O_{F_v}\arrow[r] & \Spec\mathcal O_F[1/S]
\end{tikzcd}
\]
在 \'etale 同伦型的层面是一个同伦推出（homotopy pushout）。这与 (3.2.3) 类似：那里我们遇到一个展示，其特殊之处在于它不只是推出，而是同伦推出。\footnote{为领会这一点，注意到以下事实或许有帮助：如果改用与 (3.1.1) 中不同的亏格 $g$ Riemann 曲面基本群展示，就会得到先验不同的 Betti 局部系统模空间。为什么 (3.1.1) 中的展示才是「正确」的呢？}

\end{exam}

\subsubsection{晶体条件}

下一步是切出一个在 $p$ 处施加了 $p$-进 Hodge 理论条件的导出伽罗瓦形变环。为具体起见，我们把注意力集中在「晶体」情形；它应当对应于在 $p$ 处有好约化的动机，或者换个角度说，对应于在 $p$ 处没有 level 的自守形式。

我们有映射 $\mathcal M^{\overline\rho}_{\mathcal O_F[1/S]}\to\mathcal M^{\overline\rho}_{F_p}$。（我们现在知道，对 $G=\mathrm{GL}_n$，由 [BIP23] 的工作，$\mathcal M^{\overline\rho}_{F_p}=\mathrm M^{\overline\rho}_{F_p}$ 是经典的；一般情形也预期如此，不过我们在逻辑上不会用到这一点。）我们想定义一个局部函子 $\mathcal M^{\overline\rho,\mathrm{crys}}_{F_p}\to\mathcal M^{\overline\rho}_{F_p}$，它「切出」「晶体」形变。一般情形下我们不知道如何做到这一点，因为我们无法在足够一般的环上表述晶体性的模空间论概念（§9.1 将进一步讨论）。

就目前所知，这个问题的唯一解决办法是：把限制放在那些「应当」实际上是经典的局部条件上，即「$\mathcal M^{\overline\rho,\mathrm{crys}}_{\mathcal O_F[1/S]}=\mathrm R^{\overline\rho,\mathrm{crys}}_{\mathcal O_F[1/S]}$」，其中右端是 $p$-进 Hodge 理论中某个被经典地研究的对象；然后取如下映射的复合作为我们的局部条件
\[
\mathrm R^{\overline\rho,\mathrm{crys}}_{\mathcal O_F[1/S]}\to\mathcal M^{\overline\rho}_{\mathcal O_F[1/S]}\hookrightarrow\mathcal M^{\overline\rho}_{\mathcal O_F[1/S]}.
\]

按 §3.2.1，要辨认哪些形变条件「应当」是经典的，办法是看它们是否为「正确」维数的 LCI。实践中，要分析得到的对象，还需要知道切复形；除非形变函子光滑（此时切复形就是切空间），切复形通常很难计算。

一切都能顺利运作的一类条件是 Fontaine--Laffaille 范围，[GV18] 研究的正是它。此时所有问题都消失了：
\begin{itemize}[leftmargin=1.4em]
\item 存在完整的模空间论描述，甚至是整系数的。
\item 障碍消失，于是形变函子光滑，其切复形可以显式描述。
\end{itemize}
本节余下部分都把自己置于这一情形。我们假定处在 Fontaine--Laffaille 范围内，特别地 $p$ 相对于我们所考虑的 Hodge--Tate 权是「足够大」的。如上所述，这给出闭浸入 $\mathcal M^{\overline\rho,\mathrm{crys}}_{\mathcal O_F[1/S]}\hookrightarrow\mathcal M^{\overline\rho}_{\mathcal O_F[1/S]}$。我们于是定义 $\mathcal R^{\overline\rho,\mathrm{crys}}_{\mathcal O_F[1/S]}$ 表示

\ysrc{505}下图中的导出纤维积 $\mathcal M^{\overline\rho,\mathrm{crys}}_{\mathcal O_F[1/S]}$：
\[
\begin{tikzcd}[column sep=large, row sep=large]
\mathcal M^{\overline\rho,\mathrm{crys}}_{\mathcal O_F[1/S]}\arrow[r]\arrow[d] & \mathcal M^{\overline\rho}_{\mathcal O_F[1/S]}\arrow[d]\\
\mathrm M^{\overline\rho,\mathrm{crys}}_{\mathcal O_F[1/S]}\arrow[r, hook] & \mathrm M^{\overline\rho}_{\mathcal O_F[1/S]}
\end{tikzcd}
\]
具体地说，$\mathrm R^{\overline\rho,\mathrm{crys}}_{\mathcal O_F[1/S]}$ 用导出张量积来计算（cf. §A.10）。

关键之处在于：取导出纤维积之后，我们对所得的切复形有了控制。它的 Euler 特征可用整体域伽罗瓦上同调的 Poitou--Tate 对偶来计算，答案是 $-\delta$，其中 $\delta\ge0$ 是某个称为「缺陷」（defect）的量，我们不打算详细解释它（就我们的目的而言，可以把它定义为该 Euler 特征，但文献中通常从别的出发点来定义）。量 $\delta$ 的出现至关重要；同一个 $\delta$ 也出现在「自守侧」，这是 Calegari--Geraghty 方法的枢纽。

\subsection{与 Calegari--Geraghty 方法的关系}

现在讨论 [GV18] 的主要结果。它们对 Calegari--Geraghty 方法 [CG18] 给出重新解释与扩展；后者本身是 Taylor--Wiles 方法在「$\delta>0$」情形下的修正。读者至少应熟悉 [Cal20] 层面的这些方法，才能从本小节有所收获。

\subsubsection{猜想图景}

我们简记 $\mathcal R_0:=\mathcal R^{\overline\rho,\mathrm{crys}}_{\mathcal O_F[1/S]}$，$\mathrm R_0:=\pi_0(\mathcal R_0)$。设 $\check G$ 是 $G$ 的 Langlands 对偶群。有一个复形 $\mathcal M_0$，由「$\check G$ 的 level $S$ 自守形式」组成；它由与 $\check G$ 相伴随的局部对称空间的同调组装而成。我们把 $\mathcal M_0$ 的次数安排在 $[-\delta,0]$ 中，尽管后文（§7）可能会把它重新编号到次数「$[-\delta-q_0,-q_0]$」中。关键之处在于这个复形的长度为 $\delta$——正是同一 $\delta$ 取负号出现在 $\mathcal R_0$ 的维数中。最后记 $\mathrm M_0:=H_0(\mathcal M_0)$。

Langlands 对应的「自守到伽罗瓦」方向把伽罗瓦表示附于自守形式，从而给 $\mathrm M_0$ 赋予 $\mathrm R_0$-模结构。（视 $\check G$ 是什么而定，这可能是猜想性的。）预期（大致）是：

\begin{quote}
$\mathcal R_0$ 作用在 $\mathcal M_0$ 上，并具有以下性质。
\begin{itemize}[leftmargin=1.4em]
\item 取同伦/同调群后，它恢复 $\mathrm R_0$ 在 $\mathrm M_0$ 上给定的作用，如下图：
\[
\begin{tikzcd}[column sep=small, row sep=small]
\mathcal R_0\arrow[d] & \curvearrowright & \mathcal M_0\arrow[d]\\
\mathrm R_0 & \curvearrowright & \mathrm M_0
\end{tikzcd}
\]
\item 它使 $\pi_*(\mathcal M_0)[1/p]$ 在 $\pi_*(\mathcal R_0)[1/p]$ 上一般地自由。\footnote{技术上，我们把 $\mathcal R_0$ 定义成了 pro-环，因此 $\pi_*(\mathcal R_0)$ 是 pro-群；在这样的对象上「把 $p$ 逆过来」需要一些小心。}
\end{itemize}
\end{quote}

[GV18] 的主要结果是在某些猜想与某些 Taylor--Wiles 假设下构造出 $\pi_*(\mathcal R_0)$ 在 $\pi_*(\mathcal M_0)$ 上的作用，并具有上述性质。

\ysrc{506}\begin{exam}{5.5.1}\label{thm:a6-03-11}（$\mathrm{GL}_1$ 的情形）目前，上述图景只有在 $G=\mathrm{GL}_1$ 时才能完全实现，而且即便这种情形也不平凡。\footnote{到目前为止我们关注的都是 $G$ 为伴随型的情形，因此要容纳这一情形，我们的定义须稍作修改。}若 $\delta=0$（当 $F=\Q$ 或二次虚域时如此），则所有环都是经典的，该陈述是类域论的推论。要点在于经典形变环本质上是 $\pi_1(\mathcal O_F[1/S])$ 的交换化的（完备）群环。类域论把这一交换化与 $\mathcal O_F[1/S]$ 的理想类群等同起来，而后者的 adele 描述正是相应的「局部对称空间」（一个有限点集）。

然而在其他所有情形都有 $\delta>0$，这蕴含：
\begin{itemize}[leftmargin=1.4em]
\item 导出伽罗瓦形变环不是经典的；
\item 局部对称空间不再是离散点集，
\end{itemize}
而所假设的作用超出了经典伽罗瓦形变理论。尽管如此，$\mathcal R_0$ 在 $\mathcal M_0$ 上想要的作用仍可构造出来，这将在即将发表的论文 [FHMR] 中解释。要点在于：$\mathcal R_0$ 可描述为 $\pi_1(\mathcal O_F[1/S])$ 的「导出交换化」的（完备）「导出群环」。为阐明其含义：
\begin{itemize}[leftmargin=1.4em]
\item 群的交换化控制它到交换群的态射集合。群的导出交换化控制它到动态交换群的态射构成的 anima。这可用泛性质精确表述：它是从（动态）交换群到（动态）群的具体函子的左伴随。
\item 交换群的群环也可用泛性质刻画：它是从交换环到交换群、取单位的函子的左伴随。类似地，「导出群环」构造可刻画为从动态环到动态交换群的类似函子的左伴随。
\end{itemize}
[FHMR] 的要点是：$\pi_1(\mathcal O_F[1/S])$ 的导出交换化「就是」由 $G$ 的相应局部对称空间得到的动态交换群（在 $G=\mathrm{GL}_1$ 的特殊情形，它有拓扑群结构），而其导出群环「就是」这个动态交换群上的同调链，并带有诱导的环结构。

与此同时，$\mathcal M_0$ 也正是局部对称空间上的同调链，而所求的作用就是恒等作用（tautological action）。特别地，当 $G=\mathrm{GL}_1$ 时，$\mathcal M_0$ 是 $\mathcal R_0$ 上秩一的自由模。
\end{exam}

\subsubsection{Galatius--Venkatesh 的结果}

[GV18] 主要结果的一个非正式陈述是：在关于自守提升的「标准」猜想与假设之下，并在「无同余」与 Fontaine--Laffaille 假设下，
\begin{enumerate}[leftmargin=1.8em]
\item $\pi_*(\mathcal R_0)$ 支撑在次数 $0\le *\le\delta$ 中，并且是 $\pi_1(\mathcal R_0)$ 上的外代数，而 $\pi_1(\mathcal R_0)$ 又是 $\mathrm R_0$ 上秩 $\delta$ 的自由模；
\item $H_*(\mathcal M_0)$ 带有 $\pi_*(\mathcal R_0)$ 上的自由分次模结构，并延拓 $*=0$ 时的通常结构。
\end{enumerate}

% 本 chunk 止于 5.5.2 的编号列表；该小节余下部分（自「The conjectures alluded to above...」起）见下一 chunk（印刷页 507 起）。
